% sections/02_methodology.tex
% 第2章：分析方法与比较基线

\section{分析方法与比较基线}

太空算力的可行性分析涉及多个学科——热物理、航天工程、光伏技术、电池化学、运载火箭、半导体制造——且讨论者（马斯克、黄仁勋、第三方分析师）经常使用不同的口径和假设，得出表面上互不兼容的数字。本章在进入具体分析之前，建立全文统一的分析框架和术语规范，以确保读者不会因口径混淆而在后续章节中产生误解。

\subsection{三层递进分析框架}

本报告的分析按物理 $\to$ 工程 $\to$ 商业三层递进，每一层在前一层结论成立的基础上展开：

\begin{enumerate}[leftmargin=*]
    \item \textbf{物理可行性（第4章）}：仅使用物理定律和可公开验证的物理常数做出判断。如果斯特藩-玻尔兹曼定律表明所需散热面积无法接受，或 LEO 热环境使有效散热通量远低于可工作水平，则工程和商业分析就失去前提。\emph{判断标准：物理上是否可能？}
    \item \textbf{工程可行性（第5章）}：在物理可行的前提下，分析能否从工程上实现——卫星有多重、能否装进火箭、能否在轨道上持续运行、故障后能否修复。\emph{判断标准：工程上能否造出来并持续运行？}
    \item \textbf{商业可行性（第6章）}：在工程可行的前提下，分析是否值得投入——总拥有成本（TCO）与地面数据中心的比较、哪些参数组合会改变结论方向。\emph{判断标准：商业上是否值得？}
\end{enumerate}

这一递进框架的核心价值在于避免"跨层混淆"：物理可行不代表工程可行（例如散热面积在物理上只需 0.71\;km\textsuperscript{2} 但工程上如何将 ~6,400 颗卫星的散热面精确对准深空是完全不同的问题）；工程可行不代表商业可行（例如即使每颗卫星都能造出来，总造价可能使其在经济上无法与地面竞争）。

\subsection{五大统一口径定义}

全文所有数字均基于以下统一口径。这些口径的选择理由在每个子条目中说明。

\subsubsection{功率口径}

\begin{itemize}[leftmargin=*]
    \item \textbf{GPU/计算侧名义峰值（150\;kW）}：SpaceX 公开材料中"150 kW of GPU power"所指的 GPU 直流输入端峰值功率，被广泛传播为 AI1 的核心识别数字【二档·权威资料，经 Inkl/Tom's Hardware、Knightli、Indexbox 三家独立媒体交叉验证一致】\cite{spacex2026ai1}。本文在叙事上沿用 150\;kW 这一公共讨论锚点；电源系统的工程闭合（阵列面积、PCDU 质量、光伏成本）仍按 LEO 轨道能量平衡推导，而不按 150\;kW 直接估算。

          150\;kW 与 120\;kW/210\;kW 并非同一物理量：150\;kW 是"这颗卫星计算侧能跑到多大"的名义峰值，跨越光伏阵列$\to$电池$\to$PCDU$\to$GPU 的完整能量链后，在 LEO 轨道物理中仅能支撑约 86\;kW 的持续 IT 负载。120\;kW 持续 IT 对应的正向物理需求约为 210\;kW 阵列 EOL 输出（完整推导见第5章“功率口径澄清”小节及附录 C\;\S C.2）。为直观区分这些数字，表\ref{tab:power_layers}给出完整的四层功率口径分层。
    \item \textbf{持续 IT 负载（120\;kW）}：SpaceX 官方公开的 AI1 持续算力对应的功率值\cite{spacex2026ai1}。本文以 120\;kW 作为 1\;MW 部署所需卫星数的折算分母（1\;MW / 120\;kW $\approx$ 8.33 颗）。若按 LEO 轨道能量平衡的实际约束（150\;kW 峰值仅对应 $\sim$86\;kW 持续），1\;MW 部署需约 12 颗卫星（而非 8.3 颗），对应 TCO 也会进一步上升。
    \item \textbf{IT 负载功率（1\;MW 持续，共同分母）}：最终送达计算硬件的可用功率。本报告选择 1\;MW 持续 IT 负载作为比较分母，原因有三：(1) 数字足够大以具备经济讨论意义；(2) 恰好落在 AI1（120\;kW/颗）的整数倍缩放范围内（约 8.33 颗）；(3) 地面 1\;MW 级 AI 集群已有成熟的公开 TCO 数据可作为比较参照。
\end{itemize}

\begin{table}[H]
\centering
\caption{AI1 功率口径分层——不同电气节点的物理含义}
\label{tab:power_layers}
\footnotesize
\begin{tabular}{p{3.8cm}>{\centering\arraybackslash}p{2.0cm}p{3.0cm}p{4.0cm}}
\toprule
\textbf{口径} & \textbf{数值} & \textbf{物理位置} & \textbf{报告用途} \\
\midrule
GPU/计算侧名义峰值 & 150 kW & GPU直流输入端/计算载荷侧 & 对齐SpaceX公开说法，公众记忆锚点和热峰值参考 \\
持续IT负载 & 120 kW & 实际送达计算硬件的轨道平均负载 & 1MW折算分母：1MW/120kW $\approx$ 8.33颗 \\
阵列EOL输出需求 & $\approx$210 kW & 光伏阵列/PCDU上游输出端 & 电源系统、光伏面积、PCDU质量、阵列成本计算 \\
阵列BOL名义需求 & $\approx$230 kW & 寿命初期光伏阵列输出 & 光伏成本和面积校验参考 \\
\bottomrule
\end{tabular}
\end{table}

一句话总结：150\;kW是"这颗卫星计算侧能跑到多大"的名义峰值，210\;kW是"为了长期维持120\;kW持续IT负载，光伏阵列在寿命末期需要给出的电源能力"——二者不是同一个物理量，跨越了完整能量转换链（光伏阵列$\to$电池$\to$PCDU$\to$GPU），不可混淆为同一口径。

\subsubsection{功率口径深层说明：为何保留150kW官方口径}

AI1的"150\;kW"已经成为公开传播中的核心识别数字。若本文在主叙事中立即将其替换为LEO能量平衡下的阵列需求功率（$\sim$210\;kW），容易把讨论提前带入口径争议。因此，本文在主叙事中保留150\;kW作为SpaceX公开宣称的计算侧名义峰值，并以120\;kW持续IT负载作为1\;MW折算分母；工程章节再给出完整的LEO电源闭合计算。

电源系统计算仍按LEO轨道日照/食影能量平衡闭合：若要维持120\;kW持续IT负载，计入食影供电、电池回充、转换损耗和寿命末期退化后，阵列EOL输出需求约为210\;kW。因此，本文的光伏面积、PCDU质量和光伏成本并非按150\;kW阵列输出估算，而是按物理闭合后的约210\;kW阵列需求估算。完整能量链推导见第5章“功率口径澄清”小节及附录 C\;\S C.2。

换言之，150\;kW在本文中是"公共讨论锚点"和"计算侧峰值参考"，不是电源系统的闭合功率。若将150\;kW理解为阵列峰值，则其在LEO轨道上仅能支撑约86\;kW持续IT负载，所需卫星数和TCO都会上升（误读校验表见第5章“误读校验”小节）。

\subsubsection{时间窗口径}

\begin{itemize}[leftmargin=*]
    \item \textbf{5 年分支}：单代卫星在轨服役 5 年后自动坠毁。在 10 年总窗口的比较分母下，这意味着需要发射 2 代卫星（第 1 代 + 重射 1 次）。选择 5 年的理由：LEO 卫星的大气阻力寿命、电池循环寿命和硬件退化速度在 600\;km 轨道上的工程实践表明 5 年是合理的保守估计。
    \item \textbf{10 年分支}：单代卫星在轨服役 10 年（无整星重射）。这是一个乐观但非不合理的工程目标——需要电池化学从 NMC 切换为更长寿的方案（详见第5章电池三分支分析）。
    \item \textbf{比较分母}：10 年总窗口。所有分支的 TCO 均折算到同一 10 年窗口，使得 5 年分支的重射成本被显式包含而非隐含忽略。
\end{itemize}

\subsubsection{架构代际口径}

\begin{itemize}[leftmargin=*]
    \item \textbf{主线架构：NVIDIA Blackwell（GB300 NVL72）}。这是截至 2026 年中期已公开发布且规格可独立验证的最新 NVIDIA GPU 架构。GB300 NVL72 计算托盘的详细规格来自 Lenovo 产品指南（29\;kg/托盘，4 GPU + 2 CPU）【二档·权威资料】\cite{lenovo2025gb300}，为计算载荷质量链提供了 A1 级地面硬件锚点。
    \item \textbf{不采用 Rubin 作为主线}：NVIDIA Vera Rubin 架构虽已发布（GTC 2026），但整星 BOM、空间化参数和一手公开规格尚不充分，当前仅能进行架构能效重映射（详见第5章），不能作为完整商业分析的主线。
    \item \textbf{不采用 H100 作为主线}：上一代架构，能效和算力密度显著落后于当前讨论。
\end{itemize}

\subsubsection{成本边界口径}

本报告的成本模型覆盖以下项，每一项均有完整定义和边界条件：

\begin{itemize}[leftmargin=*]
    \item \textbf{制造成本}：含 9 个子系统——计算载荷、散热系统、光伏阵列、电池、功率电子/配电（PCDU）、平台/结构、推进、通信/激光链路、系统裕量/冗余/附加（AIT裕量：涵盖冗余硬件、系统质量裕量与集成附加质量）。每项均按证据档位标注。
    \item \textbf{发射成本}：\$200/kg 基准（Starship V3 复用目标），敏感性分析覆盖 \$100--500/kg。此值的选择依据见第6章“发射单价敏感性分析”与本章“发射单价口径说明”。
    \item \textbf{保险成本}：发射成本的 8\%（基于航天保险行业惯例）。
    \item \textbf{运维成本}：\$5M/年（地面测控、星座管理和软件更新）。
    \item \textbf{退役成本}：\$3M/代（受控再入）。
    \item \textbf{不包含}：地面网关/用户段、频谱许可费、研发非经常性工程费用（NRE）——这些或因高度场景依赖、或因缺乏公开数据而无法纳入统一的比较框架。
\end{itemize}

\subsubsection{发射单价口径}

\$200/kg 被用作基准发射单价。这一选择基于以下三层证据：（1）历史层——发射成本在过去 15 年已下降约 100 倍（从 Shuttle 约 \$54,500/kg 到 Falcon 9 复用约 \$2,720/kg）\cite{nasalaunch2018}；（2）经济阈值层——\textless\$200/kg 被广泛认为是使大规模 LEO 部署进入经济可行区的重要阈值；（3）工程目标层——Starship V3 的 "100+ 吨全复用"目标对应约 \$100--200/kg 的边际成本{}\cite{33fg2026starship}。本报告将 \$200/kg 定位为"成熟复用工程目标"而非"当前已实现值"，并在第6章的敏感性分析中覆盖更宽泛的 \$100--500/kg 区间。

\subsection{口径对齐矩阵（简版）}

表\ref{tab:alignment_brief}给出了本报告主链与主要外部口径的简化对比。不同口径下"同一个 MW"的物理含义截然不同——如果两个分析使用不同的时间窗、功率定义或架构代际，直接比较其 TCO 数字在方法上无效。

\begin{figure}[H]
\centering
% 口径对齐示意图 — 第3章用
% 展示不同口径（时间窗、功率、架构、成本）之间的对应关系与不可直接比较的警示
% 性冷淡风：灰/石板蓝/深岩灰配色
\begin{tikzpicture}[node distance=0.8cm, auto]
  % 定义样式
  \tikzstyle{block} = [rectangle, draw=gray!60, fill=cyan!18!gray!8, text width=2.35cm, align=center, rounded corners, minimum height=2.55cm, font=\small, text=black!85]
  \tikzstyle{arrow} = [thick,draw=gray!50,->,>=stealth]
  \tikzstyle{warn}  = [font=\small\bfseries, text=black!60]

  % 四个主要口径块
  \node [block] (time) {\textbf{时间窗}\\[4pt]5年 vs 10年\\含/不含重射};
  \node [block, right=of time] (power) {\textbf{功率口径}\\[4pt]峰值150 kW\\持续120 kW\\IT负载1 MW};
  \node [block, right=of power] (arch) {\textbf{架构代际}\\[4pt]Blackwell 基线\\Rubin 对照};
  \node [block, right=of arch] (cost) {\textbf{成本边界}\\[4pt]制造/发射/\\运维/退役};

  % 连接箭头
  \draw [arrow] (time) -- (power);
  \draw [arrow] (power) -- (arch);
  \draw [arrow] (arch) -- (cost);

  % 底部标注
  \node [text width=10.6cm, align=center, font=\small, anchor=north] at ($(time.south)!0.5!(cost.south)+(0,-1.0cm)$) {
    \textbf{\textcolor{black!70}{不同口径假设下的数字不可直接比较。}}\\[4pt]
    本报告所有比较均基于统一基线：\\
    Blackwell架构、10年窗口、1 MW持续IT、\$200/kg发射。
  };
\end{tikzpicture}

\caption{\protect\parbox{0.92\textwidth}{\centering 口径对齐示意图——不同口径假设下的数字不可直接比较，本报告所有比较均基于统一基线（Blackwell架构、10年窗口、1 MW持续IT、\$200/kg发射）。}}
\label{fig:alignment_chart}
\end{figure}

\begin{table}[H]
\centering
\caption{口径对齐矩阵（简版）——各口径在关键维度上的取值与可比较性}
\label{tab:alignment_brief}
\footnotesize
\begin{tabular}{p{2.0cm}p{1.8cm}p{1.8cm}p{1.8cm}p{1.8cm}p{1.8cm}}
\toprule
\textbf{维度} & \textbf{本报告主链} & \textbf{马斯克公开口径} & \textbf{黄仁勋公开口径} & \textbf{此前行业分析} & \textbf{常见舆论口径} \\
\midrule
时间窗 & 10 年 & 未明确 & "数年" & 10 年 & 5 年 \\
功率口径 & 1\;MW 持续 IT & GW 级部署量 & 未量化 & 1\;MW IT & 1\;MW \\
架构代际 & Blackwell & 未指定 & Rubin & Blackwell & 未指定 \\
光伏路线 & GaAs 主链 & AI1（未公开） & 未公开 & GaAs & HJT \\
发射成本 & \$200/kg & 乐观 & 未论述 & \$200/kg & \$50--100/kg \\
寿命假设 & 10年单代 & 未明确 & 保留 & 10 年 & 5 年 \\
可比较性 & 基准 & 换算后可比 & 部分可比较 & 部分可比较 & 换算后可比 \\
\bottomrule
\end{tabular}
\end{table}

完整版口径对齐矩阵（7 个场景 $\times$ 9 个维度，含来源标注和定量影响分析）见附录 A。

\subsection{口径差异为何重要}

在阅读后续章节之前，有必要阐明一个核心方法学观点：本报告中不同数字之间的矛盾，绝大多数不是计算错误，而是口径差异。例如：

\begin{itemize}[leftmargin=*]
    \item 一个分析得出"10 年总成本 \$765M"，另一个得出"\$1,341M"——差异主要来自前者假设 10 年寿命、后者假设 5 年寿命（需要重射一次）。两者都可能是"正确的"，但比较时需要显式声明寿命假设。
    \item 一个分析说"发射成本可以降到 \$50/kg"，另一个说"\$200/kg"——前者可能隐含数千次发射的学习曲线效应，后者可能只考虑初期 69 次/年的 FAA 批准上限。两者讨论的是不同时间点的不同情景。
    \item 一个分析说"HJT 比 GaAs 便宜 70--100 倍"，另一个说"GaAs 是唯一可行路线"——前者从电池组件成本出发，后者从整星翼展和部署约束出发。两者在各自的局部范围内都有道理，但全局最优解需综合考虑所有约束。
\end{itemize}

本报告的处理方式是将分歧分解为可独立参数化的维度（寿命、电池、电源路线、质量链、价格链），逐维度核查，并给出每个维度的参数可以独立变化时结论如何改变。这一方法使分析的"分歧空间"从模糊的立场对立收缩为可量化的参数敏感性区间。

\subsection{方法论声明：第一性原理自底向上验证}

本报告所采用的分析方法——从物理定律（斯特藩-玻尔兹曼）、工程硬件锚点（NVIDIA GB300 NVL72官方规格）、到独立Python建模的完整计算链——遵循第一性原理自底向上验证的方法论。这一方法论确保：

\begin{enumerate}[leftmargin=*]
    \item \textbf{主结论尽量追溯到公开可独立核验的原始来源}——大多数关键结论可沿附录D的五层路由从正文结论追溯到原始URL或物理公式；少数链路（约30\%的四档推断参数）仍依赖内部裁决记录与待开源脚本。
    \item \textbf{即使四档参数（14个合理推断参数）的取值发生变动，也不会引起数量级层面的结论翻转}——因为 TCO 的成本结构中，光伏阵列（$\sim$50\%）已通过二档/三档交叉锚定；AIT裕量仍是主导成本项（$\sim$20\%），但其取值仍带有四档推断属性。
    \item \textbf{与基于简单比例外推或线性缩放的"通用估算"方法不同}，本报告的估算在每个环节都显式标注了证据档位、推断依据和失效风险。
    \item \textbf{本报告追求的不是"精确到小数点"的预测精度，而是"即使参数在合理区间内变化，方向性结论不翻转"的稳健性。}
\end{enumerate}

\endinput
